\scanpage{15}
\[
L_{DE}(t)=L_E\bigl(p^{-(n\unsure{-|+}\rho)}t\bigr)
\]
d'où, portant \unsure{dans (2)|du cas} et remplaçant $t$ par $t^{-1}$ :
\[
\boxed{L_E\left(\frac{1}{p^{n+\rho}t}\right)=\varepsilon(E)\left(-p^{\frac{n+\rho}{2}}t\right)^{\chi(E)}L_E(t)}\tag*{(3)}
\]

Posant
\[
\boxed{\xi_E(t)=t^{\chi(E)/2}L_E(t)}
\]
il vient l'équation fonctionnelle
\[
\boxed{\xi\left(\frac{1}{p^{n+\rho}t}\right)=\varepsilon'(E)\,\xi(t)}
\]
\[
\begin{aligned}
\varepsilon'(E)=\pm1,\qquad \varepsilon'(E)&=(-1)^{\chi(E)}\prod_i\operatorname{signe}\operatorname{dét} f_{H^i(\bar X,\bar F)}\\
&=(-1)^{\chi(E)}\operatorname{signe}\operatorname{det} f_{H^*(\bar X,\bar F)}
\end{aligned}
\]

\deleted{Si on \unsure{veut} l'équation fonctionnelle avec $\varepsilon(E)$ au lieu de} Notons d'ailleurs que $H^i(\bar X,\bar E)$ et $H^{2n-i}(\bar X,\check{\bar E})$ \inserted{$\simeq H^{2n-i}(\bar X,\bar E(\rho))=H^{2n-i}(\bar X,\bar E)(\rho)$} étant ``\unsure{duals|duaux|dual} \unsure{à|:} valeurs \unsure{dans|au} $\mathbb{Q}_{\ell}(n)$'', et l'opération de $f$ sur \inserted{$\mathbb{Q}_{\ell}(\unsure{1|n})$} \unsure{--|en|étant} de signe positif, on trouve que $\varepsilon_i(E)=\varepsilon_{2n-i}(E)$, et par suite
\[
\boxed{\varepsilon(E)=\operatorname{Signe}\operatorname{det} f_{H^n(\bar X,\bar E)}}.
\]
Pour la raison analogue, $(-1)^{\chi(E)}=(-1)^{b_n}$. Donc
\[
\varepsilon'(E)=(-1)^{b_n(E)}\operatorname{signe}\operatorname{dét}f_{H^n(X,E)}
\]
% archive sheet number: 14
